\chapter{How selection and reproduction change covariance}
\label{inheritance:chapter}

Selection changes which parents contribute to reproduction. To calculate the
genetic variation among their offspring, we also need to specify what they
transmit. These are separate biological steps. A distribution of parental
breeding values weighted by expected reproductive output describes the first
step; segregation, recombination and mating determine the second.

This chapter follows covariance through both steps under a stated inheritance
model. The central question is whether variation can persist when stabilising
selection repeatedly reduces covariance among contributing parents. We first
identify the genetic variance that the model holds fixed, then derive the
update. Matrix order gives a positive covariance equilibrium without requiring
the trait axes to be independent. For one trait, a stronger argument gives a
unique equilibrium and convergence from every nonnegative starting variance.
The boundary with no reduction of associations after selection shows why
these conclusions depend on reproduction.

\section{What is held fixed when allelic effects are fixed?}
\label{inheritance:genic}

Let the breeding value be an $n$-component column vector. In the additive
genetic architecture considered here, a copy of a designated allele at locus
$j$ contributes the fixed vector $\vect\alpha_j$. If $N_j\in\{0,1,2\}$ is
its diploid allele count, then, up to a constant vector,
\begin{equation}
  \vect A=\sum_{j=1}^{m}\vect\alpha_jN_j,
  \qquad
  \vect X=\vect A-\E[\vect A]
     =\sum_{j=1}^{m}\vect\alpha_j(N_j-2p_j),
  \label{inheritance:additive-map}
\end{equation}
where both homologous allele copies have marginal frequency $p_j$.
The fixed-effect assumption makes this the same linear map in each
generation. It is an extra assumption about genetic architecture: a
population-specific regression definition of breeding value does not, by
itself, guarantee fixed allelic coefficients across generations.

Write $N_j=Z_{j1}+Z_{j2}$, where $Z_{jh}$ is a zero--one indicator for the
allele on homologue $h$. Under random association of the two homologues,
\begin{align}
  \E[N_j]&=2p_j,\\
  \Var(N_j)
    &=\Var(Z_{j1})+\Var(Z_{j2})
       +2\Cov(Z_{j1},Z_{j2})\\
    &=2p_j(1-p_j).
\end{align}
If allele counts at different loci are also uncorrelated, the covariance
of the sum in Equation~\eqref{inheritance:additive-map} is
\begin{equation}
  \mat B=2\sum_{j=1}^{m}p_j(1-p_j)
                       \vect\alpha_j\vect\alpha_j^{\mathsf T}.
  \label{inheritance:genic-formula}
\end{equation}
This is the genic covariance baseline used in the model. Each outer product
$\vect\alpha_j\vect\alpha_j^{\mathsf T}$ describes how one allele's effects
contribute to variances and covariances of traits; the factor
$2p_j(1-p_j)$ describes variation in its copy number. An allele with effects
on several traits can therefore contribute off-diagonal genic covariance
even when different loci are independent.

Actual allele counts need not have these random-association properties.
Expanding the covariance without assuming independence gives
\begin{align}
  \mat G
    &=\sum_j\Var(N_j)\vect\alpha_j\vect\alpha_j^{\mathsf T}
      +\sum_{j<k}\Cov(N_j,N_k)
       \bigl(\vect\alpha_j\vect\alpha_k^{\mathsf T}
            +\vect\alpha_k\vect\alpha_j^{\mathsf T}\bigr),\\
  \mat G&=\mat B+\mat D.
  \label{inheritance:association-decomposition}
\end{align}
Here $\mat D$ is the total departure from the genic baseline. It includes
the cross-locus terms and, if homologues are associated within loci, the
terms $2\Cov(Z_{j1},Z_{j2})\vect\alpha_j\vect\alpha_j^{\mathsf T}$.
Unlike a covariance matrix, $\mat D$ need not be positive semidefinite:
associations can reduce variation in a trait combination. This distinction
underlies the effect of selection-induced disequilibrium on quantitative
genetic variance \citep{Bulmer1971,Bulmer1974}.

\begin{example}[Fixed effects with a changing genic variance]
  \label{inheritance:frequency-example}
  Consider one trait and one locus with allelic effect $\alpha=1$.
  At $p=0.5$, the genic variance is $2(0.5)(0.5)=0.50$.
  At $p=0.8$, it is $2(0.8)(0.2)=0.32$.
  The effect is unchanged, but its variance contribution is smaller.
\end{example}

\paragraph{Biological interpretation.}
Holding the allelic effects fixed keeps breeding-value coordinates
comparable across generations. Holding $\mat B$ fixed additionally removes
changes in genic covariance caused by changing allele frequencies. The
latter can represent an approximation over a short interval, or a formal
infinitesimal model in which genic covariance remains constant. It is not a
consequence of additive gene action alone. Mutation, drift, migration and
changes in allelic effects would require further terms or a different model
\citep{BartonTurelli1991,WalshLynch2018}.

\section{From selected parents to the next generation}
\label{inheritance:closure}

At generation $t$, assume that breeding values before selection are
Gaussian, with mean $\vect\mu_t$ and positive-definite covariance
$\mat G_t$. Write $\vect x=\vect a-\vect\mu_t$. In this chapter the
log conditional mean of reproductive output is exactly quadratic on the
population's support:
\begin{equation}
  \ell_A(\vect\mu_t+\vect x)
     =\ell_{0,t}+\vect b_t^{\mathsf T}\vect x
                      -\frac12\vect x^{\mathsf T}\bGamma\vect x,
  \qquad \bGamma\succeq\mat0.
  \label{inheritance:quadratic-fitness}
\end{equation}
Thus $\bGamma=-\mat H$ is the sign-reversed Hessian and describes stabilising
curvature. The exact assumption matters: substituting a Taylor expansion
does not make the selected distribution exactly Gaussian.

Weighting the Gaussian density by
$W_A(\vect a)=\exp\{\ell_A(\vect a)\}$ adds the quadratic precision term
$\bGamma$. The resulting covariance among parents, under these reproductive
weights, is
\begin{equation}
  \mat G_t^{\mathrm{sel}}
       =(\mat G_t^{-1}+\bGamma)^{-1}.
  \label{inheritance:selected-covariance}
\end{equation}
The precision is positive definite because $\mat G_t^{-1}\succ\mat0$ and
$\bGamma\succeq\mat0$. Consequently the expected-output weights have a
finite normalising integral for every finite $\vect b_t$. Completing the
square changes the mean, but the coefficient of the centred quadratic
determines covariance and contains no $\vect b_t$.

Equation~\eqref{inheritance:selected-covariance} has not yet described
offspring. To do so, hold the genic covariance at $\mat B\succ\mat0$ and
make the following assumption about the departure from this baseline:
\begin{equation}
  \mat D_{t+1}=\rho\mat D_t^{\mathrm{sel}},
  \qquad
  \mat D_t^{\mathrm{sel}}=\mat G_t^{\mathrm{sel}}-\mat B,
  \qquad 0\leq\rho<1.
  \label{inheritance:retention-assumption}
\end{equation}
The selected departure includes associations already present as well as
their change during the current episode of selection. The assumption retains
the same fraction $\rho$ of this whole departure after reproduction.
Substitution into $\mat G_{t+1}=\mat B+\mat D_{t+1}$ gives
\begin{align}
  \mat G_{t+1}
    &=\mat B+\rho(\mat G_t^{\mathrm{sel}}-\mat B)\\
    &=(1-\rho)\mat B+
                    \rho(\mat G_t^{-1}+\bGamma)^{-1}.
  \label{inheritance:recursion}
\end{align}
This derivation identifies what the inheritance model contributes. The
Gaussian selection calculation supplies $\mat G_t^{\mathrm{sel}}$;
Equation~\eqref{inheritance:retention-assumption} specifies how reproduction
uses it. Neither a variance share nor a fitness regression supplies that
transmission assumption.

\subsection{Why one half is a useful benchmark}
\label{inheritance:half-retention}

A moment calculation explains the familiar freely recombining,
random-mating infinitesimal benchmark \citep{Bulmer1971}. Suppose two
contributing parents are independent draws from the same selected
distribution. In a model for their offspring, write
\begin{equation}
  \vect A_{\mathrm{o}}
      =\frac{\vect A_1+\vect A_2}{2}+\vect M,
  \qquad
  \E[\vect M\mid\text{parents}]=\vect0,
  \label{inheritance:mendelian-decomposition}
\end{equation}
where $\vect M$ is the Mendelian segregation deviation. The conditional
zero mean makes its covariance with the midparent value zero. If its
mean conditional covariance is $\mat B/2$, the law of total covariance
gives
\begin{align}
  \Cov(\vect A_{\mathrm{o}})
     &=\Cov\!\left(\frac{\vect A_1+\vect A_2}{2}\right)
          +\E[\Cov(\vect M\mid\text{parents})]\\
     &=\frac14\bigl(\mat G_t^{\mathrm{sel}}
                         +\mat G_t^{\mathrm{sel}}\bigr)+\frac12\mat B\\
     &=\frac12\mat G_t^{\mathrm{sel}}+\frac12\mat B.
  \label{inheritance:half-reproduction}
\end{align}
This is Equation~\eqref{inheritance:recursion} with $\rho=1/2$.
The two ingredients---independent parents and a specified segregation
covariance---are part of the model. For example, at an independently
segregating locus a heterozygous parent transmits the allele with variance
$1/4$ in copy number. Averaging over two parents with heterozygote frequency
$2p_j(1-p_j)$ gives $p_j(1-p_j)$, half the locus's genic copy-number
variance. Summing the corresponding effect outer products gives
$\mat B/2$ under independent segregation across loci and the fixed
single-locus heterozygosities assumed by this infinitesimal closure.

The more general scalar $\rho$ in Equation~\eqref{inheritance:recursion}
summarises retention; it is not derived from a single recombination fraction.
Different locus pairs have different recombination rates and contributions
to covariance. Nonrandom mating can also change the covariance between
parents and the segregation term. These mechanisms do not, in general,
rescale every covariance departure by the same number
\citep{Bulmer1974,BartonTurelli1991}. In particular, this model supplies no
universal conversion from a selfing rate or an inbreeding coefficient to
$\rho$.

\paragraph{Biological interpretation.}
When selected associations reduce variance, reproduction can restore some
variance by reducing those associations. The term $(1-\rho)\mat B$ does not
represent new mutations or new alleles entering the population. It expresses
the assumed return towards the covariance that the same allele frequencies
and effects would produce under random association. Maintaining Gaussian
breeding-value distributions in successive generations is also part of this
closure; a rule for covariance alone does not prove that distributional
assumption.

\section{How to compare covariance matrices}
\label{inheritance:matrix-order}

For one trait, a smaller variance has an immediate meaning. With several
traits, comparing individual matrix entries is insufficient: a covariance
matrix may lose variance in one trait combination and gain it in another.
The useful ordering compares every linear combination at once; it is the
positive-semidefinite order used in matrix and convex analysis
\citep{BoydVandenberghe2004}.

\begin{definition}[Loewner order]
  \label{inheritance:loewner}
  For symmetric matrices $\mat C$ and $\mat D$, write
  $\mat C\preceq\mat D$ when $\mat D-\mat C$ is positive semidefinite,
  that is,
  \begin{equation}
    \vect v^{\mathsf T}\mat C\vect v
       \leq\vect v^{\mathsf T}\mat D\vect v
       \quad\text{for every }\vect v\in\R^n.
  \end{equation}
  Strict positive definiteness is written $\mat C\prec\mat D$.
\end{definition}

Because $\Var(\vect v^{\mathsf T}\vect A)=
\vect v^{\mathsf T}\mat G\vect v$, the order says that every specified
linear combination of traits has at least as much variance under
$\mat D$ as under $\mat C$. It is a partial order: two covariances need
not be comparable. The same comparison is preserved by a change of
coordinates. Indeed, if $\mat D-\mat C\succeq\mat0$, then
\begin{equation}
  \vect v^{\mathsf T}\mat R(\mat D-\mat C)\mat R^{\mathsf T}\vect v
   =(\mat R^{\mathsf T}\vect v)^{\mathsf T}
       (\mat D-\mat C)(\mat R^{\mathsf T}\vect v)\geq0.
  \label{inheritance:congruence-order}
\end{equation}
Multiplication on the two sides in this form is called a congruence.
It does not require the matrices to commute.

\begin{lemma}[Stabilising selection reduces covariance]
  \label{inheritance:inverse-order}
  If $\mat G\succ\mat0$ and $\bGamma\succeq\mat0$, then
  \begin{equation}
    \mat0\prec (\mat G^{-1}+\bGamma)^{-1}\preceq\mat G.
    \label{inheritance:selection-order}
  \end{equation}
  If $\bGamma\succ\mat0$, the upper inequality is strict.
  Equality of the two complete covariance matrices occurs exactly when
  $\bGamma=\mat0$.
\end{lemma}

\begin{proof}
  A positive-definite symmetric matrix has an orthogonal eigenvector basis
  with positive eigenvalues. Taking their positive square roots defines
  $\mat G^{1/2}$ and its inverse $\mat G^{-1/2}$. Define
  \begin{equation}
    \mat K=\mat G^{1/2}\bGamma\mat G^{1/2}\succeq\mat0.
  \end{equation}
  Factor the precision in the order shown:
  \begin{align}
    \mat G^{-1}+\bGamma
       &=\mat G^{-1/2}(\mat I+\mat K)\mat G^{-1/2},\\
    (\mat G^{-1}+\bGamma)^{-1}
       &=\mat G^{1/2}(\mat I+\mat K)^{-1}\mat G^{1/2}.
    \label{inheritance:whitened-selection}
  \end{align}
  The eigenvalues $\kappa_i$ of $\mat K$ are nonnegative. Therefore the
  eigenvalues of $(\mat I+\mat K)^{-1}$ are
  $1/(1+\kappa_i)$, which lie in $(0,1]$. It follows that
  \begin{equation}
    \mat0\prec(\mat I+\mat K)^{-1}\preceq\mat I.
  \end{equation}
  Congruence by $\mat G^{1/2}$ gives
  Equation~\eqref{inheritance:selection-order}. If $\bGamma\succ\mat0$,
  every $\kappa_i>0$, so every eigenvalue $1/(1+\kappa_i)<1$ and the
  reduction is positive definite. Equality of the full matrices requires
  every $\kappa_i=0$. Since $\mat G^{1/2}$ is invertible, this is
  equivalent to $\bGamma=\mat0$.
\end{proof}

This proof is an instance of the reversal of Loewner order by inversion:
if $\mat P\succeq\mat Q\succ\mat0$, then
$\mat P^{-1}\preceq\mat Q^{-1}$. To see the general statement, apply the
same argument to
$\mat Q^{-1/2}\mat P\mat Q^{-1/2}\succeq\mat I$ and undo the
congruence. The proof diagonalises one symmetric matrix; it never assumes
that $\mat G$ and $\bGamma$ share eigenvectors.

\paragraph{Biological interpretation.}
Under the Gaussian quadratic model, stabilising selection cannot increase
variance in any trait combination among contributing parents. A reduction in one
episode of selection does not imply a reduction from one generation to the
next: reproduction subsequently changes the associations. Nor does zero
curvature along one coordinate guarantee an unchanged marginal variance
there, because that trait can covary with a trait under selection.

\section{A positive matrix equilibrium exists}
\label{inheritance:matrix-fixed-point}

Define the covariance update on positive-definite matrices by
\begin{equation}
  F(\mat G)=(1-\rho)\mat B+
                       \rho(\mat G^{-1}+\bGamma)^{-1}.
  \label{inheritance:map}
\end{equation}
Assume throughout this section that $\mat B\succ\mat0$,
$\bGamma\succeq\mat0$ and $0\leq\rho<1$ are fixed.
The lower matrix
\begin{equation}
  \mat L=(1-\rho)\mat B
\end{equation}
is positive definite. Consider the set
\begin{equation}
  \mathcal K=\{\mat G=\mat G^{\mathsf T}:\mat L\preceq\mat G\preceq\mat B\}.
  \label{inheritance:interval}
\end{equation}
The two bounds concern every trait combination, not entrywise intervals.

\begin{proposition}[Existence and bounds]
  \label{inheritance:existence}
  The map $F$ has at least one positive-definite fixed point in
  $\mathcal K$. Every positive-definite fixed point of $F$, whether or
  not considered initially within $\mathcal K$, satisfies
  \begin{equation}
    (1-\rho)\mat B\preceq\mat G^*\preceq\mat B.
    \label{inheritance:fixed-bounds}
  \end{equation}
\end{proposition}

\begin{proof}
  We first show that $F$ leaves the set $\mathcal K$ invariant. For
  $\mat G\in\mathcal K$, the selected covariance is positive definite
  and no larger than $\mat G$ by Lemma~\ref{inheritance:inverse-order}.
  Hence
  \begin{align}
    F(\mat G)&\succeq(1-\rho)\mat B=\mat L,\\
    F(\mat G)&\preceq(1-\rho)\mat B+\rho\mat G\\
             &\preceq(1-\rho)\mat B+\rho\mat B=\mat B.
    \label{inheritance:invariance}
  \end{align}
  Thus $F(\mathcal K)\subseteq\mathcal K$.

  Next, regard real symmetric $n\times n$ matrices as a Euclidean space
  with $n(n+1)/2$ coordinates, one for each independent matrix entry.
  The set $\mathcal K$ is nonempty because it contains $\mat B$.
  It is convex: for $\mat G_1,\mat G_2\in\mathcal K$ and
  $0\leq\theta\leq1$, each of
  \begin{align}
    \theta\mat G_1+(1-\theta)\mat G_2-\mat L
      &=\theta(\mat G_1-\mat L)+(1-\theta)(\mat G_2-\mat L),\\
    \mat B-\theta\mat G_1-(1-\theta)\mat G_2
      &=\theta(\mat B-\mat G_1)+(1-\theta)(\mat B-\mat G_2)
  \end{align}
  is a sum of positive-semidefinite matrices.

  It is also closed: a limit of matrices satisfying each quadratic-form
  inequality in Equation~\eqref{inheritance:interval} still satisfies it.
  It is bounded because $\mat0\preceq\mat G\preceq\mat B$ implies
  $\lambda_{\max}(\mat G)\leq\lambda_{\max}(\mat B)$ and therefore
  \begin{equation}
    \norm{\mat G}_{\mathrm F}^2
       =\sum_{i=1}^{n}\lambda_i(\mat G)^2
       \leq n\lambda_{\max}(\mat B)^2.
  \end{equation}
  Closed and bounded subsets of a finite-dimensional Euclidean space
  are compact. In this setting, compactness means that any sequence in
  the set has a convergent subsequence whose limit stays in the set.

  Finally, matrix inversion is continuous wherever the determinant is
  nonzero. Throughout $\mathcal K$,
  \begin{equation}
    \lambda_{\min}(\mat G)
       \geq(1-\rho)\lambda_{\min}(\mat B)>0,
  \end{equation}
  so the first inversion in $F$ is well-defined. The second is also
  well-defined because $\mat G^{-1}+\bGamma\succ\mat0$.
  Matrix sums and scalar multiplication are continuous, so $F$ is
  continuous on $\mathcal K$.

  Brouwer's fixed-point theorem states that a continuous map from a
  nonempty compact convex subset of a finite-dimensional Euclidean
  space into itself has a fixed point \citep{GranasDugundji2003}.
  All its hypotheses have now been checked. The resulting
  $\mat G^*\in\mathcal K$ is positive
  definite because it is no smaller than $\mat L\succ\mat0$.

  To bound any other positive-definite fixed point, no assumption
  about its starting position is needed. The equation
  $\mat G^*=F(\mat G^*)$ immediately gives
  $\mat G^*\succeq(1-\rho)\mat B$. Applying
  Equation~\eqref{inheritance:selection-order} at $\mat G^*$ gives
  \begin{align}
    \mat G^*&\preceq(1-\rho)\mat B+\rho\mat G^*,\\
    (1-\rho)\mat G^*&\preceq(1-\rho)\mat B,\\
    \mat G^*&\preceq\mat B.
  \end{align}
  Division by $1-\rho$ preserves the inequality because this scalar
  is strictly positive. This proves both bounds for every positive
  fixed point.
\end{proof}

\paragraph{Biological interpretation.}
For any trait combination $\vect v^{\mathsf T}\vect A$, an equilibrium
variance lies between $(1-\rho)\vect v^{\mathsf T}\mat B\vect v$ and
$\vect v^{\mathsf T}\mat B\vect v$. Thus fixed positive genic covariance
and incomplete retention prevent a covariance equilibrium from losing all
variation in a genetic direction. If $\rho>0$, its lower bound is strict
because the selected covariance is positive definite. If $\rho=0$, the
next covariance is exactly $\mat B$. With $\rho>0$, the full matrix
$\mat G^*=\mat B$ is possible exactly when $\bGamma=\mat0$.

The argument proves existence and bounds. Brouwer's theorem does not say
that the fixed point is unique or that repeated updates converge to it.
Invariance says that a trajectory starting in $\mathcal K$ stays there;
it does not say that every other starting covariance enters that interval
after one update. We make no general matrix uniqueness or convergence
claim here. The scalar case permits a separate, stronger proof.

\section{One trait: the fixed point and every trajectory}
\label{inheritance:scalar}

Let $g$ be a breeding-value variance, $B>0$ the genic variance, and
$\gamma\geq0$ the stabilising curvature. The scalar map is
\begin{equation}
  F(g)=(1-\rho)B+\rho\frac{g}{1+\gamma g},
  \qquad g\geq0,\quad0\leq\rho<1.
  \label{inheritance:scalar-map}
\end{equation}
It also describes an isotropic model with
$\mat G=g\mat I$, $\mat B=B\mat I$ and
$\bGamma=\gamma\mat I$ in fixed trait coordinates.
The product $\gamma g$ is dimensionless: $g$ and $B$ have units of
trait squared, while $\gamma$ has inverse squared-trait units.
Although the precision calculation used $g>0$, the fraction in
Equation~\eqref{inheritance:scalar-map} has the continuous value zero
at $g=0$. The scalar recursion therefore includes the zero starting
variance, followed by $F(0)=(1-\rho)B>0$.

\subsection{Solving the fixed-point equation}
\label{inheritance:scalar-root}

Set $g^*=F(g^*)$ and first take $\gamma>0$. Since
$1+\gamma g^*>0$, multiplying by the denominator preserves the equation:
\begin{align}
  g^*&=(1-\rho)B+\rho\frac{g^*}{1+\gamma g^*},\\
  g^*+\gamma(g^*)^2
     &=(1-\rho)B+(1-\rho)\gamma B g^*+\rho g^*,\\
  \gamma(g^*)^2+(1-\rho)(1-\gamma B)g^*-(1-\rho)B&=0.
  \label{inheritance:scalar-quadratic}
\end{align}
The leading coefficient is positive and the constant term is negative.
Thus the product of the two roots is negative, so exactly one root is
positive. Write $c=(1-\rho)(1-\gamma B)$ and
$d=(1-\rho)B>0$. The positive root is
\begin{equation}
  g^*=\frac{-c+\sqrt{c^2+4\gamma d}}{2\gamma}.
  \label{inheritance:positive-root}
\end{equation}
When $c\geq0$, subtracting two nearly equal positive numbers in this
formula can lose floating-point precision. Multiplying numerator and
denominator by $\sqrt{c^2+4\gamma d}+c$ gives the equivalent form
\begin{equation}
  g^*=\frac{2d}{\sqrt{c^2+4\gamma d}+c}.
  \label{inheritance:rationalised-root}
\end{equation}
Use the rationalised form when $c\geq0$ and the first form when $c<0$
to avoid that subtraction. This is an arithmetic choice, not a different
equilibrium. As $\gamma$ tends to zero with $\rho<1$ fixed, the
rationalised expression tends to $B$, agreeing with the direct
zero-curvature calculation below.

\subsection{Why the scalar iteration converges}
\label{inheritance:contraction}

An equilibrium equation alone does not tell us where a trajectory goes.
To answer that question, compare two possible current variances $g,h\geq0$:
\begin{align}
  F(g)-F(h)
    &=\rho\left(\frac{g}{1+\gamma g}
                      -\frac{h}{1+\gamma h}\right)\\
    &=\rho\frac{g(1+\gamma h)-h(1+\gamma g)}
                    {(1+\gamma g)(1+\gamma h)}\\
    &=\frac{\rho(g-h)}{(1+\gamma g)(1+\gamma h)}.
\end{align}
The two mixed terms cancel. Since each denominator factor is at least
one,
\begin{equation}
  |F(g)-F(h)|\leq\rho|g-h|.
  \label{inheritance:contraction-bound}
\end{equation}
A map that reduces every distance by at most a constant strictly below
one is a contraction.

\begin{proposition}[Global scalar convergence]
  \label{inheritance:scalar-convergence}
  For $B>0$, $\gamma\geq0$ and $0\leq\rho<1$, the map in
  Equation~\eqref{inheritance:scalar-map} has a unique fixed point on
  $[0,\infty)$. Every nonnegative starting variance converges to it,
  with error bound
  \begin{equation}
    |g_t-g^*|\leq\rho^t|g_0-g^*|,\qquad t\geq1.
    \label{inheritance:scalar-error-bound}
  \end{equation}
\end{proposition}

\begin{proof}
  The map takes $[0,\infty)$ into itself. This set is complete in
  the usual distance $|g-h|$: a Cauchy sequence of nonnegative real
  numbers has a nonnegative real limit. These properties and
  Equation~\eqref{inheritance:contraction-bound} are the hypotheses
  of the Banach contraction theorem \citep{Banach1922,GranasDugundji2003}.
  Its argument can be seen directly in the iterates.

  For $0<\rho<1$, repeated use of the bound gives
  \begin{equation}
    |g_{k+1}-g_k|\leq\rho^k|g_1-g_0|.
  \end{equation}
  For integers $p\geq1$, the triangle inequality then gives
  \begin{align}
    |g_{t+p}-g_t|
       &\leq\sum_{k=t}^{t+p-1}|g_{k+1}-g_k|\\
       &\leq|g_1-g_0|\sum_{k=t}^{t+p-1}\rho^k\\
       &\leq\frac{\rho^t}{1-\rho}|g_1-g_0|.
  \end{align}
  The final quantity tends to zero as $t\to\infty$, independently
  of $p$. The iterates are therefore Cauchy and converge to a value
  $g^*\geq0$. Continuity gives
  $F(g^*)=\lim_tF(g_t)=\lim_tg_{t+1}=g^*$.

  If two fixed points $g^*$ and $h^*$ existed, then
  \begin{equation}
    |g^*-h^*|=|F(g^*)-F(h^*)|\leq\rho|g^*-h^*|.
  \end{equation}
  As $1-\rho>0$, their distance must be zero. Finally, compare each
  iterate with this fixed point in
  Equation~\eqref{inheritance:contraction-bound} and apply the
  inequality $t$ times to obtain
  Equation~\eqref{inheritance:scalar-error-bound}. If $\rho=0$,
  $F(g)=B$ for every $g$, so all conclusions follow after one update.
\end{proof}

The convergence proof needs neither a positive starting variance nor
a starting point inside $[(1-\rho)B,B]$. It also explains the direction
of change. For $0<\rho<1$, the exact difference formula gives
\begin{equation}
  F(g)-g^*
     =\frac{\rho(g-g^*)}{(1+\gamma g)(1+\gamma g^*)}.
  \label{inheritance:no-overshoot}
\end{equation}
The sign of the deviation from equilibrium is preserved and its magnitude
strictly decreases unless it is already zero. Thus a start below $g^*$
increases monotonically towards it; a start above $g^*$ decreases
monotonically towards it. Neither trajectory crosses the fixed point.
At $\rho=0$, it reaches the fixed point immediately.

\paragraph{Biological interpretation.}
Stabilising selection never increases the selected-parent variance in this
model, but reproduction need not leave the next generation below its
preselection value. With $B=1$, $\gamma=0.3$, $\rho=0.9$ and
$g_t=0.1$, selection gives $0.1/1.03$, and reproduction gives
\begin{equation}
  g_{t+1}=0.1+0.9\frac{0.1}{1.03}
           =0.1873786408.
  \label{inheritance:increase-example}
\end{equation}
The starting variance is below the equilibrium. Convergence describes
both increasing and decreasing trajectories; it does not imply that
covariance must fall in every population.

\subsection{Boundary cases and the worked iteration}
\label{inheritance:scalar-boundaries}

With no retention, $\rho=0$, reproduction returns the variance to $B$
after one step. With no curvature, $\gamma=0$ and $\rho<1$,
\begin{align}
  g_{t+1}&=(1-\rho)B+\rho g_t,\\
  g_{t+1}-B&=\rho(g_t-B),\\
  g_t&=B+\rho^t(g_0-B).
  \label{inheritance:zero-curvature}
\end{align}
Therefore $B$ is the unique fixed point, even if the original population
contains covariance departures that reproduction gradually removes.
The matrix version of this particular boundary is equally direct:
$\mat G_t=\mat B+\rho^t(\mat G_0-\mat B)$ when
$\bGamma=\mat0$. This explicit boundary calculation is separate from
any general matrix convergence claim.

For the worked example in the mathematics guide, $B=1$,
$\gamma=0.5$, $\rho=0.5$ and $g_0=1$. The first selected variance is
$2/3$, so the first offspring variance is
$g_1=1/2+(1/2)(2/3)=5/6$. Further substitutions give
\begin{center}
\begin{tabular}{ccc}
  \toprule
  Generation & Exact variance & Decimal variance\\
  \midrule
  0 & $1$ & $1.000000$\\
  1 & $5/6$ & $0.833333$\\
  2 & $27/34$ & $0.794118$\\
  3 & $149/190$ & $0.784211$\\
  4 & $827/1058$ & $0.781664$\\
  \bottomrule
\end{tabular}
\end{center}
The fixed-point polynomial reduces to
$2(g^*)^2+g^*-2=0$, giving
\begin{equation}
  g^*=\frac{\sqrt{17}-1}{4}=0.7807764064.
  \label{inheritance:guide-fixed-point}
\end{equation}
These iterates remain between $0.5$ and $1$ because that interval is
invariant and contains the starting value. The numerical sequence
illustrates the proven convergence; it is not its proof.

\section{Retaining every covariance departure}
\label{inheritance:rho-one}

The boundary $\rho=1$ removes the baseline term entirely. It must be
analysed separately: the positive lower bound $(1-\rho)\mat B$ and the
strict scalar contraction constant used above are then lost. Assume
$\mat G_0\succ\mat0$ and fixed $\bGamma\succeq\mat0$. The update becomes
\begin{equation}
  \mat G_{t+1}=(\mat G_t^{-1}+\bGamma)^{-1}.
\end{equation}
Inverting both sides gives the simple precision recursion
\begin{equation}
  \mat G_{t+1}^{-1}=\mat G_t^{-1}+\bGamma.
\end{equation}
The initial case is $\mat G_0^{-1}=\mat G_0^{-1}+0\bGamma$.
If $\mat G_t^{-1}=\mat G_0^{-1}+t\bGamma$, one further update gives
$\mat G_{t+1}^{-1}=\mat G_0^{-1}+(t+1)\bGamma$.
Induction therefore proves
\begin{equation}
  \mat G_t=(\mat G_0^{-1}+t\bGamma)^{-1}.
  \label{inheritance:precision-iteration}
\end{equation}
No simultaneous diagonalisation of $\mat G_0$ and $\bGamma$ was used.

If $\bGamma\succ\mat0$, then for $t>0$
\begin{equation}
  \mat0\prec\mat G_t\preceq(t\bGamma)^{-1}
                     =\frac1t\bGamma^{-1}.
  \label{inheritance:positive-curvature-limit}
\end{equation}
The right-hand side tends to zero, so covariance disappears in every
direction. If $\bGamma=\mat0$, every covariance stays at its starting
value. In one trait with $\gamma>0$,
$g_t=g_0/(1+t\gamma g_0)$ for a positive start. The example
$g_0=1$, $\gamma=0.5$ gives $g_4=1/3$. At this boundary the continuous
scalar map also leaves $g_0=0$ at zero.

\subsection{Singular curvature: what variation actually remains?}
\label{inheritance:singular-limit}

Suppose $\bGamma$ has rank $r$ with $0<r<n$. Let the columns of
$\mat U$ be orthonormal eigenvectors with positive curvature, and
let the columns of $\mat V$ be an orthonormal basis for its null space.
Then $\mat Q=[\mat U\;\mat V]$ is orthogonal and
\begin{equation}
  \bGamma=\mat U\mat\Lambda\mat U^{\mathsf T},
  \qquad \mat\Lambda\succ\mat0,
  \qquad
  \mat Q^{\mathsf T}\mat G_0\mat Q
     =\begin{pmatrix}\mat C&\mat J\\\mat J^{\mathsf T}&\mat D\end{pmatrix}.
  \label{inheritance:curvature-blocks}
\end{equation}
Here $\mat C$ is the initial covariance of the coordinates with positive
curvature, $\mat D$ that of the zero-curvature coordinates, and $\mat J$
their cross-covariance. The block letters in this subsection are local
matrix names; $\mat D$ here is not the association departure used earlier.

\begin{proposition}[The exact retained covariance]
  \label{inheritance:schur-limit}
  Under these assumptions,
  \begin{equation}
    \lim_{t\to\infty}\mat G_t
       =\mat V\bigl(\mat D-\mat J^{\mathsf T}\mat C^{-1}\mat J\bigr)
                        \mat V^{\mathsf T}.
    \label{inheritance:null-limit}
  \end{equation}
  The matrix in parentheses is positive definite on the null-space
  coordinates. It equals their initial Gaussian conditional covariance
  given the positive-curvature coordinates, rather than their marginal
  covariance $\mat D$ in general.
\end{proposition}

\begin{proof}
  For $t>0$ define
  \begin{equation}
    \mat M_t=t^{-1}\mat\Lambda^{-1}+\mat U^{\mathsf T}\mat G_0\mat U
            =t^{-1}\mat\Lambda^{-1}+\mat C.
  \end{equation}
  Both terms are positive definite, so this matrix can be inverted.
  We will use the identity
  \begin{equation}
    (\mat G_0^{-1}+t\mat U\mat\Lambda\mat U^{\mathsf T})^{-1}
       =\mat G_0-\mat G_0\mat U\mat M_t^{-1}\mat U^{\mathsf T}\mat G_0.
    \label{inheritance:inverse-update}
  \end{equation}
  This is a form of the matrix inverse update identity, often called
  the Woodbury identity. Here it can be verified without invoking a
  general formula. Multiplying its right-hand side on the left by
  $\mat G_0^{-1}+t\mat U\mat\Lambda\mat U^{\mathsf T}$ gives
  \begin{equation}
    \mat I+\mat U
       \bigl[t\mat\Lambda-\mat M_t^{-1}
                   -t\mat\Lambda\mat C\mat M_t^{-1}\bigr]
         \mat U^{\mathsf T}\mat G_0.
  \end{equation}
  The bracket vanishes because
  \begin{equation}
    t\mat\Lambda\mat M_t=\mat I+t\mat\Lambda\mat C,
    \qquad
    (\mat I+t\mat\Lambda\mat C)\mat M_t^{-1}=t\mat\Lambda.
  \end{equation}
  This establishes Equation~\eqref{inheritance:inverse-update}, with
  matrix products in their stated order.

  As $t\to\infty$, $\mat M_t\to\mat C\succ\mat0$, so
  $\mat M_t^{-1}\to\mat C^{-1}$. Expressing the limiting identity
  in the $\mat Q$ coordinates gives
  \begin{align}
    \lim_{t\to\infty}\mat Q^{\mathsf T}\mat G_t\mat Q
       &=\begin{pmatrix}\mat C&\mat J\\\mat J^{\mathsf T}&\mat D\end{pmatrix}
         -\begin{pmatrix}\mat C\\\mat J^{\mathsf T}\end{pmatrix}
             \mat C^{-1}\begin{pmatrix}\mat C&\mat J\end{pmatrix}\\
       &=\begin{pmatrix}
            \mat0&\mat0\\
            \mat0&\mat D-\mat J^{\mathsf T}\mat C^{-1}\mat J
          \end{pmatrix}.
  \end{align}
  Transforming back gives Equation~\eqref{inheritance:null-limit}.

  The matrix $\mat S=\mat D-\mat J^{\mathsf T}\mat C^{-1}\mat J$
  is the Schur complement of $\mat C$. For a nonzero null-coordinate
  vector $\vect z$, use the positive definiteness of the full initial
  covariance to obtain
  \begin{equation}
    \vect z^{\mathsf T}\mat S\vect z
      =\begin{pmatrix}-\mat C^{-1}\mat J\vect z\\\vect z\end{pmatrix}^{\mathsf T}
        \begin{pmatrix}\mat C&\mat J\\\mat J^{\mathsf T}&\mat D\end{pmatrix}
        \begin{pmatrix}-\mat C^{-1}\mat J\vect z\\\vect z\end{pmatrix}>0.
  \end{equation}
  Finally, write $\vect X_{\mathrm{curved}}=\mat U^{\mathsf T}\vect X_0$
  and $\vect X_{\mathrm{null}}=\mat V^{\mathsf T}\vect X_0$ for the
  centred initial coordinates. In a joint Gaussian distribution,
  subtracting the conditional linear mean of the null coordinates leaves residual
  $\vect X_{\mathrm{null}}-\mat J^{\mathsf T}\mat C^{-1}
     \vect X_{\mathrm{curved}}$.
  Its covariance is $\mat S$, and its covariance with
  $\vect X_{\mathrm{curved}}$ is zero. Joint Gaussianity then makes the
  residual independent of those coordinates, proving the conditional
  covariance interpretation.
\end{proof}

\begin{example}[A trait with no curvature can still lose variance]
  \label{inheritance:singular-example}
  Let
  \begin{equation}
    \mat G_0=\begin{pmatrix}1&0.4\\0.4&2\end{pmatrix},
    \qquad
    \bGamma=\begin{pmatrix}0.5&0\\0&0\end{pmatrix}.
  \end{equation}
  The second trait has zero curvature but initially covaries with the
  first. Its retained variance is
  $2-0.4^2/1=1.84$, so
  \begin{equation}
    \lim_{t\to\infty}\mat G_t
          =\begin{pmatrix}0&0\\0&1.84\end{pmatrix}.
  \end{equation}
  Keeping the original marginal variance $2$ would ignore the initial
  association with the trait whose variation is removed. When
  $\mat J=\mat0$, marginal and conditional covariance do coincide.
\end{example}

\paragraph{Biological interpretation.}
With no restoration, persistent covariance is confined to trait
combinations with zero curvature. Its amount depends on the original
associations with combinations under stabilising selection. This is a
covariance limit under the specified model, not a statement that allele
frequencies, mean traits or reproductive output reach a joint equilibrium.
If the initial covariance is singular, first restrict the calculation to
its positive-variance support. At $\rho=1$ no baseline term adds variation
outside that support; use the curvature restricted to the same coordinates
before taking inverses.

\section{A covariance trajectory does not determine the fitness gradient}
\label{inheritance:gradient-and-share}

The recursion contains $\mat G_t$, $\mat B$, $\bGamma$ and $\rho$, but no
directional gradient $\vect b_t$. This independence follows from the
exact Gaussian and quadratic selection calculation in
Equation~\eqref{inheritance:selected-covariance}. It is not a general
property of selection on non-Gaussian populations or nonquadratic fitness
functions.

The first-order variance share still depends on the gradient. In the
isotropic case with $n$ traits,
\begin{align}
  \Vlin(t)&=g_t\norm{\vect b_t}^{2},\\
  \Vquad(t)&=\frac12 n\gamma^2g_t^2,\\
  \Ag(t)&=\frac{\norm{\vect b_t}^{2}}
                   {\norm{\vect b_t}^{2}+\tfrac12n\gamma^2g_t},
                   \qquad g_t>0,
  \label{inheritance:isotropic-share}
\end{align}
where the denominator must be positive. If $\vect b_t=\vect0$ but
$\gamma>0$, the share is zero. If both terms vanish, the variance share
is undefined. Under fixed nonzero gradient and positive curvature, a
decrease in $g_t$ increases the share even though both component variances
decrease. A start below the scalar equilibrium instead produces increasing
$g_t$ and a decreasing share under that same fixed gradient.

Figure 4 of the paper chooses $n=4$, $g_0=4$, $B=1$,
$\gamma=0.30$ and $\rho=0.90$. The scalar equilibrium is
$g^*=0.4723532349$. To calculate the share, it separately specifies
\begin{equation}
  \norm{\vect b_t}=
  \begin{cases}
    1,&t\leq40,\\
    \exp[-0.06(t-40)],&t>40.
  \end{cases}
  \label{inheritance:figure-gradient-history}
\end{equation}
While the gradient is maintained, the share rises from $0.581395$ to
$0.921639$ by generation 40. It subsequently falls as the imposed
gradient decreases, reaching approximately $0.00870449$ at generation
100. The figure combines a derived covariance trajectory with a supplied
gradient history. That history is not an output of the covariance recursion.

A changing local gradient need not mean that the fitness surface changes.
For a fixed quadratic function
$\ell_A(\vect a)=c+\vect d^{\mathsf T}\vect a
-\tfrac12\vect a^{\mathsf T}\bGamma\vect a$, differentiation gives
\begin{equation}
  \vect b_t=\vect d-\bGamma\vect\mu_t.
  \label{inheritance:gradient-at-mean}
\end{equation}
Movement of the population mean changes the local gradient on the same
surface. Predicting that movement would require its own transmission
assumptions and mean recursion. The covariance equilibrium proved here
does not provide them.

\section{Checks and the scope of the conclusions}
\label{inheritance:checks}

The repository script
\codepath{code/reproduce_result_checks.py} records the group
\codepath{inheritance_matrix_and_limits}. It checks the
single-locus variances $0.50$ and $0.32$, twelve matrix interval cases with
noncommuting baseline and curvature, one numerical fixed point and its
bounds, scalar contractions and several starting values, and the correlated
singular-curvature limit $1.84$. The chosen matrices are
\begin{equation}
  \mat B=\begin{pmatrix}1&0.25\\0.25&0.6\end{pmatrix},
  \qquad
  \bGamma=\begin{pmatrix}0.6&-0.15\\-0.15&0.3\end{pmatrix},
  \qquad \rho=0.65.
\end{equation}
They are positive definite and do not commute. Thus their numerical check
does not rely on independent trait directions. It complements the proof
for all matrices satisfying the assumptions; a finite list of checked
cases could not establish that general result.

The source \codepath{code/figures.py} generates the main
trajectory, while \codepath{code/guide_figures.py} generates
the scalar worked iteration. The main reproduction script also checks
archived forward-time population simulations. Those simulations compare
a current-population Gaussian variance formula with empirical moments.
Their agreement does not independently validate the fixed-$\mat B$
inheritance closure derived here.

The covariance results therefore answer a conditional biological question:
what happens when a fixed genic baseline, stabilising curvature and a
specified retention rule act on Gaussian breeding-value distributions?
Within that model, incomplete retention allows a positive covariance
equilibrium, and the scalar trajectory converges to it from either side.
To ask about changing allele frequencies, linkage architecture or mating
systems, those mechanisms must enter the model explicitly. The derivations
identify the assumptions that connect selection among parents to variation
in the next generation.
